\paragraph{Operacje punktowe.} Wstawienie i~zmiana pojedynczego rekordu
po kluczu (C1, C2, C4--C6, U1, U2, U4, U5, J3) we wszystkich bazach trwają
poniżej 2\,ms i~nie zależą od rozmiaru zbioru. Pojedyncza rezerwacja (C1) na
10 mln biletów trwa \val{valkey}{10000000}{off}{C1}\,ms w~Valkey,
\val{postgres}{10000000}{off}{C1}\,ms w~PostgreSQL,
\val{mariadb}{10000000}{off}{C1}\,ms w~MariaDB i~\val{mongo}{10000000}{off}{C1}\,ms
w~MongoDB. MongoDB jest tu najwolniejszy, bo każdy zapis czeka na
potwierdzenie z~dziennika (\texttt{majority}), a~C1 dodatkowo jest
transakcją wielodokumentową. Valkey nie zapisuje niczego na dysk, więc jego
przewaga w~zapisach wynika częściowo z~innego poziomu trwałości.

\paragraph{Zapytania bez indeksów rosną liniowo.} W~bazach relacyjnych
i~w~MongoDB każde zapytanie po kolumnie innej niż klucz główny wymaga pełnego
skanu. Wolne miejsca na wydarzenie (R2) w~PostgreSQL rosną z~\val{postgres}{500000}{off}{R2}\,ms
(500 tys.) do \val{postgres}{10000000}{off}{R2}\,ms (10 mln). W~MariaDB rosną
z~\val{mariadb}{500000}{off}{R2} do \val{mariadb}{10000000}{off}{R2}\,ms,
a~w~MongoDB z~\val{mongo}{500000}{off}{R2} do \val{mongo}{10000000}{off}{R2}\,ms.
W~Valkey ten sam scenariusz trwa stale około \val{valkey}{10000000}{off}{R2}\,ms,
bo zbiór wolnych miejsc wydarzenia jest podstawową strukturą modelu, a~nie
indeksem.

\paragraph{MariaDB bez indeksów jest najwolniejsza z~relacyjnych.} Przy 10 mln
biletów MariaDB wykonuje R1 w~\val{mariadb}{10000000}{off}{R1}\,ms, a~PostgreSQL
w~\val{postgres}{10000000}{off}{R1}\,ms. Złączenie rezerwacji z~biletami bez
indeksu MariaDB realizuje pętlą zagnieżdżoną z~buforem bloków (\texttt{ALL} na
tabeli \texttt{bilety}, ok. 10 mln przeczytanych wierszy, tab.~\ref{tab:plans-rel}).
PostgreSQL używa złączenia haszowego i~równoległego skanu. Ta sama różnica
widać w~R5, U3, U6 i~D1--D6.

\paragraph{Osadzanie zamiast złączeń.} Odczyt rezerwacji z~biletami (R1)
w~MongoDB trwa \val{mongo}{10000000}{off}{R1}\,ms, a~w~Valkey
\val{valkey}{10000000}{off}{R1}\,ms, także bez indeksów drugorzędnych.
Bilety są osadzone w~dokumencie albo w~JSON-ie rezerwacji, więc wystarcza
wyszukiwanie po kluczu głównym. Ten sam efekt widać przy anulowaniu
rezerwacji (D1): MongoDB zwalnia bilety po identyfikatorach z~osadzonej
migawki, przez indeks \texttt{\_id}, w~\val{mongo}{10000000}{off}{D1}\,ms.
PostgreSQL bez indeksu na \texttt{bilety.rezerwacja\_id} potrzebuje do tego
\val{postgres}{10000000}{off}{D1}\,ms, a~MariaDB \val{mariadb}{10000000}{off}{D1}\,ms.
To przykład kompromisu normalizacji i~denormalizacji: odczyt agregatu jest
tani, ale dane biletów są przechowywane w~dwóch miejscach.

\paragraph{Agregacje.} Ranking 10 najpopularniejszych wydarzeń (R6) wymaga
w~bazach relacyjnych i~w~MongoDB przejścia przez wszystkie sprzedane bilety.
PostgreSQL z~indeksem robi to w~\val{postgres}{10000000}{on}{R6}\,ms (Index
Only Scan), MariaDB w~\val{mariadb}{10000000}{on}{R6}\,ms, MongoDB
w~\val{mongo}{10000000}{on}{R6}\,ms. Valkey odczytuje gotowy licznik ze zbioru
posortowanego w~\val{valkey}{10000000}{on}{R6}\,ms. Utrzymuje go przy każdej
sprzedaży, czyli przenosi koszt agregacji na zapis.

\paragraph{Valkey bez struktur indeksujących.} Zapytania po atrybucie
(historia klienta R4, wygaszanie U6, usuwanie D3, D4, D6, agregacja J2)
wymagają w~Valkey przejścia \texttt{SCAN} przez wszystkie klucze rezerwacji
i~zdekodowania każdego JSON-a. Przy 10 mln biletów to około
\val{valkey}{10000000}{off}{R4}\,ms na operację, najgorszy wynik w~całym badaniu.
Po dodaniu zbiorów \texttt{ix:*} ta sama historia klienta trwa
\val{valkey}{10000000}{on}{R4}\,ms. W~modelu klucz--wartość brak ,,indeksu''
nie oznacza wolniejszego planu, tylko brak jakiejkolwiek ścieżki dostępu.

\paragraph{Zmienność.} Mediany, odchylenia standardowe i~p95 są w~pliku
\texttt{generated/data/scenarios\_summary.csv}. Najwyższy współczynnik
zmienności mają operacje punktowe trwające ułamek milisekundy. Pojedyncze
powtórzenie trafiające na zimną pamięć podręczną albo punkt kontrolny WAL
wydłuża się wtedy kilkukrotnie. Dla operacji trwających powyżej 10\,ms
rozrzut jest mały i~nie zmienia kolejności baz.
